\documentclass{article}
\usepackage[margin=2cm]{geometry}
\usepackage{amsmath}
\usepackage{graphicx}
\usepackage{booktabs}
\usepackage[utf8]{inputenc}
\usepackage{ragged2e}
\setlength{\emergencystretch}{3em}
\begin{document}
\sloppy

\noindent\hfill\includegraphics[height=1.8em]{assets/figure_header_logo.png}\par\noindent\hrule\vspace{4pt}

result given by


\begin{equation}
\begin{array} { r l } { Z = \frac { 1 } { | \mathcal { W } | } \int _ { \mathrm { C a r t o n } } [ d \sigma ] ~ e ^ { - \frac { \beta \mathbf { g } _ { \mathrm { F M } } ^ { 2 } I } { \mathrm { T r } } ( \sigma ^ { 2 } ) } \mathrm { d e t } _ { \mathrm { A d } } \left ( \sin ( \mathrm { i } \pi \sigma ) \mathrm { e } ^ { \frac { 1 } { 4 } \mathrm { i } ( \sigma ) } \right ) } & { } \\ { \times \prod _ { i } \mathrm { d e t } _ { \mathrm { R}}}\end{array}
\end{equation}


the radius of. $S ^ { \mathrm { i } }$ $\alpha$ is a dimensionless matrix, and $f$is de


\begin{equation}
f ( y ) = { \frac { i \pi y ^ { 3 } } { 3 } } + y ^ { 2 } \log \left ( 1 - e ^ { - 2 \pi i y } \right ) + { \frac { i y } { \pi } } \mathrm { L i } _ { 2 } \left ( e ^ { - 2 \pi i y } \right ) + { \frac { 1 } { 2 \pi ^ { 2 } } } \mathrm { L i } _ { 3 } \left ( e ^ { - 2 \pi i y } \right ) - { \frac { \zeta ( 3 ) } { 2 \pi ^ { 2 } } } \qquad \quad ( 2 )
\end{equation}


The quotient by the Weyl group in amounts to division by a simple numerical factor. $| \mathcal { W } | =$
$2 ^ { n } N!$.The integral over is not restricted to a Weyl chamber. Though this localization.$\sigma$
result was obtained in the IR theory, it is expected to hold in the UV due to the assumed
$Q . { \ell }$-exactness of the irrelevant UV completion terms. One may rewrite the partition function.
in terms of the free energy as


\begin{equation}
\begin{array} { c } { Z = \frac { 1 } { | \mathcal { W } | } \int _ { \mathrm { C a r t a n } } [ \mathrm { d } \sigma ] ~ e ^ { - F ( \sigma ) } + O \left ( e ^ { \frac { - 4 \theta s ^ { 3 } x } { g _ { Y M } ^ { 2 } } } \right ) } \\ { F ( \sigma ) = \frac { 4 \pi ^ { 3 } r } { g _ { Y M } ^ { 2 } } \mathrm { T r } \sigma ^ { 2 } + \mathrm { T r } _ { \Lambda } F _ { V } ( \sigma ) + \sum _ { I } \mathrm {}}\end{array}
\end{equation}


The definitions of. $F . { 1 ^ { \prime } } ( \sigma )$ and $F . { I i } ( \sigma )$ follow simply from , and using one may obtain the
following large argument expansions


\begin{equation}
F _ { V } ( \sigma ) \approx \frac { \pi } { 6 } | \sigma | ^ { 3 } - \pi | \sigma | \qquad \qquad F _ { H } ( \sigma ) \approx - \frac { \pi } { 6 } | \sigma | ^ { 3 } - \frac { \pi } { 8 } | \sigma | \qquad \qquad ( 4 )
\end{equation}


It was argued in that in the large. $s . { i } ^ { i }$ limit, the perturbative Yang-Mills termi.e. the first
term in the expression for $F ( \sigma )$ incan be neglected, as can be the instanton contributions.
Thus in our evaluation of the free energy, we will only concern ourselves with the contributions coming from $F . { \mathbf { F } ^ { \prime } } ( \sigma )$ and $F . { \mathbf { H } } ( \boldsymbol { \sigma } )$ The first step in the evaluation of is recasting the
matrix integral in a simpler form. The integral over$\alpha$ in is an integration over the Coulomb
branch, which is parameterized by the non-zero vevs of . One may write$\sigma$


\begin{equation}
\begin{array} { r c c } { \sigma = \mathrm { d i a g } \{ \lambda _ { 1 }, \dots, \lambda _ { N }, - \lambda _ { 1 }, \dots, - \lambda _ { N } \} } & { \qquad } & { \qquad } & { ( 5 ) } \\ {.. \qquad.. \qquad } & { \qquad.. \qquad } & { \cdots \quad } \end{array}
\end{equation}


since $U S p ( 2 N )$ has$N$ elements in its Cartan. The integration variables are these. $i \textbf { \lambda } . { i }$$N$
Normalizing the weights of the fundamental representation of.$U S p ( 2 N )$ to be$\pm \mathbf { e } . { i }$with $e . { i }$
forming a basis of unit vectors for , it follows that the adjoint representation has weights$\mathbb { R } ^ { N }$
$\pm 2 e . { i }$ and$\pmb { e } . { i } \pm \pmb { e } . { j }$ for all$i \neq j$, whereas the anti-symmetric representation has only weights

\end{document}
